\section{Related Work and Study Positioning}

This work is positioned between three related research streams: static Android exposure measurement, runtime network-behavior analysis, and privacy-oriented app ecosystem studies. The study is also a follow-on to two ScytaleDroid predecessor papers. The first predecessor used static analysis to assess six Android social media apps, focusing on permissions, manifest configuration, exported components, strings/resources, network/storage flags, MASVS mapping, and CVSS-style severity framing \cite{day2025staticSocial}. The second predecessor studied dynamic network behavior in 12 Android social media and messaging apps using non-root physical-device captures, idle and interactive sessions, time-window traffic features, Isolation Forest and one-class SVM modeling, and a Runtime Deviation Index \cite{day2026dynamicRdi}. Paper 3 should therefore be read as an expansion and bridge, not as an exact reproduction of either predecessor.

Recent Android security work supports the need for static inspection but also shows why static evidence is incomplete by itself. ManiScope demonstrated that manifest misconfigurations remain common across Google Play and pre-installed applications \cite{yang2022maniscope}. Wang et al.\ showed that runtime-permission handling remains error-prone across modern Android versions \cite{wang2023runtimepermissions}. These studies motivate the static side of the present work: permissions, exported components, manifest attributes, network-security settings, and packaged resources are useful indicators of potential exposure, but they do not fully describe app behavior during execution.

Dynamic analysis literature emphasizes the complementary limitations of runtime evidence. Sutter et al.\ identify code coverage, stimulation strategy, environmental dependence, nondeterminism, and provenance as recurring challenges in Android dynamic security analysis \cite{sutter2024dynamic}. These constraints directly shape the present methodology: runtime captures are interpreted as scenario-bounded behavioral evidence rather than as exhaustive behavior models. The non-root physical-device setting further reflects a tradeoff between realism and observability. It avoids invasive instrumentation and decrypted content capture, but it limits visibility into application internals.

Recent privacy studies also support the need to link static app artifacts to runtime or ecosystem context. Paci et al.\ demonstrate that third-party tracking behavior in mobile apps requires execution-time observation \cite{paci2023tracking}. Rodriguez et al.\ examine privacy settings in Android third-party library integrations, showing that SDK configuration and defaults affect privacy outcomes at scale \cite{rodriguez2025privacysettings}. Khedkar et al.\ study privacy-related data collected by Android apps and emphasize the difficulty of accurately reporting app data collection practices \cite{khedkar2026privacydata}. These works motivate the present paper's privacy boundary: the study reports transport-level and evidence-bundle provenance without claiming decrypted payload semantics.

The closest recent study to Paper 3 is the 2026 empirical comparison of Android messaging apps by Karyotakis et al., which combines static and dynamic analysis under reproducible scenarios for Messenger, Signal, and Telegram \cite{karyotakis2026messaging}. Paper 3 differs in scope and governance. It expands beyond messaging to a 15-app consumer cohort, includes news and social/content apps, and focuses on preserving build-backed cutoff evidence so that app-update churn does not invalidate the paper's claims. The contribution is therefore not only the combined static/runtime measurement, but also the evidence-tier model that separates operational current-build status from paper-usable build/version-backed evidence.
